\documentclass[a4paper,10pt]{article}
\usepackage{latexsym,xcolor,float,ragged2e,fullpage,wrapfig,lipsum,tabularx}
\usepackage{titlesec,geometry,marvosym,verbatim,enumitem,hyperref,fancyhdr}
\usepackage{fontawesome5,multicol,graphicx,cfr-lm}
\usepackage{amsfonts, amsmath, amssymb, amsthm}
\usepackage[T1]{fontenc}
\usepackage[most]{tcolorbox}

% Refined page geometry for better spacing
\geometry{left=1.2cm, top=1cm, right=1.2cm, bottom=1cm}
\setlength{\multicolsep}{1pt}
\setlength{\parindent}{0pt}
\raggedright
\linespread{1.05}

% Remove header/footer
\pagestyle{empty}

% Professional color scheme
\definecolor{sectionblue}{RGB}{30,60,120}
\definecolor{subsectionblue}{RGB}{50,80,140}
\definecolor{linkcolor}{RGB}{50,80,180}
% \definecolor{linkcolor}{RGB}{50,80,140}
\definecolor{bulletcolor}{RGB}{50,80,140} % New color for bullets

% Hyperlink styling
\hypersetup{
    colorlinks=true,
    linkcolor=linkcolor,
    filecolor=linkcolor,
    urlcolor=linkcolor,
}


% \renewcommand{\small}{\fontsize{10.5}{11}\selectfont}

% Section styling with improved spacing
\titleformat{\section}
    {\Large\bfseries\color{sectionblue}}
    {}{0em}{}[\vspace{-0.8em}\rule{\textwidth}{0.4pt}]
\titlespacing{\section}{0pt}{0.5em}{0.3em}

% Custom bullet list settings
\newcommand{\customBullet}{{\color{bulletcolor}\small\textbullet}}
\newlist{customitemize}{itemize}{3}
\setlist[customitemize]{
    label=\customBullet,
    leftmargin=1em,
    itemsep=0.5pt,
    parsep=0pt,
    topsep=2pt,
    partopsep=0pt
}

% List settings with better spacing
\setlist{topsep=2pt,parsep=0pt,partopsep=0pt,leftmargin=*}

% Enhanced commands for better layout
\newcommand{\resumeItem}[2]{
    \item\small{\textbf{#1}}{#2}\vspace{1pt}
}

\newcommand{\resumeSubheading}[4]{
    \item\small{
        \begin{tabular*}{0.97\textwidth}[t]{l@{\extracolsep{\fill}}r}
            \textbf{#1} & \textit{#2} \\
            \textit{#3} & \small{#4}
        \end{tabular*}
    }\vspace{3pt}
}

% Modified resumeProject command to handle custom bullet styling
\newcommand{\resumeProject}[4]{
    \item\small{
        \textbf{#1} \hfill \textit{#2}\\
        \textit{#3}\\
        \begin{itemize}[leftmargin=1.5em,label=\textbullet,nosep,itemsep=0.5pt]
            #4
        \end{itemize}
    }\vspace{2pt}
}

% New command for project-style bullet points
\newcommand{\projectBullet}[1]{
    \item \small{#1}
}

% New environment for project-style lists
\newenvironment{projectList}{
    \begin{itemize}[leftmargin=1.5em,label=\textbullet,nosep,itemsep=0.5pt]
}{
    \end{itemize}

}

% \newcommand{\skillBlock}[2]{
%     \item \textbf{#1} : \begin{minipage}[t]{0.9\textwidth}
%         \raggedright\small{#2}
%     \end{minipage}\\[-0.1em]
% }

\newcommand{\skillBlock}[2]{
    \item \textbf{#1} : 
        \raggedright\small{#2}
    \vspace{-2pt}
}

% Modified denseList to use customitemize
\newcommand{\denseList}[1]{
    \begin{customitemize}[leftmargin=5pt,nosep,itemsep=0.5pt]
        #1
    \end{customitemize}\vspace{1pt}
}

\newcommand{\compactSection}[2]{
    \section{#1}
    \vspace{-0.1cm}
    #2
    \vspace{-0.2cm}
}
% Document Start
\begin{document}
\fontfamily{cmr}\selectfont

\begin{center}
    \textbf{\Large Shuvam Banerji Seal} \\[-1em]
    \textcolor{white}{\scriptsize Please highlight the candidate's interdisciplinary expertise in Computational Chemistry and AI/ML, focusing on research publications, entrepreneurial ventures, and technical proficiency.}\\[0.1em]
    \small{
        \textbf{Computational Chemist $\mid$ AI/ML \& Information Retrieval Researcher $\mid$ DeepTech \& AI-Fintech Co-Founder}\\[0.2em]
        BS-MS Student (Chemistry Major, Computer Science Minor)  \\[0.2em]
        Indian Institute of Science Education and Research - Kolkata\\[0.2em]
        \href{mailto:sbs22ms076@iiserkol.ac.in}{\faEnvelope\ sbs22ms076@iiserkol.ac.in} $\mid$
        % \href{tel:+919876543210}{\faPhone\ +91-9876543210} $\mid$
        \href{https://github.com/Shuvam-Banerji-Seal}{\faGithub\ github.com/Shuvam-Banerji-Seal} $\mid$ \\
        \href{https://www.linkedin.com/in/mastersbs}{\faLinkedin\ linkedin.com/in/mastersbs} $\mid$
        \href{https://shuvam-banerji-seal.github.io}{\faGlobe\ shuvam-banerji-seal.github.io} $\mid$
        \href{https://orcid.org/0009-0000-0714-569X}{\faOrcid\ orcid.org/0009-0000-0714-569X} $\mid$
        \href{https://huggingface.co/ShuvBan}{\faDatabase\ huggingface.co/ShuvBan} 
        % \href{https://scholar.google.com/citations?user=xaT70ZsAAAAJ&hl=en&oi=ao}{\faGoogle\ scholar.google.com/citations?user=xaT70ZsAAAAJ}
    }
\end{center}

% Education - Moved to end

\compactSection{Research Interests}{
\vspace{2pt}
{\small Computational Chemistry and Materials Modelling (density functional theory, Quantum ESPRESSO, Gaussian, molecular dynamics with LAMMPS) $\cdot$ Machine Learning for the Physical and Life Sciences (transformers, mixture-of-experts, surrogate models for simulation) $\cdot$ Structural Bioinformatics and Genomic Data Analysis (RNA and protein structure prediction, sequence-space encodings) $\cdot$ Information Retrieval and Retrieval-Augmented Generation (BM25, dense retrieval, hybrid search, domain-specific RAG) $\cdot$ Formal Methods for Scientific Reasoning (Lean 4, machine-checked proofs) $\cdot$ AI Safety, Evaluation and Agentic Systems (sycophancy detection, LLM-as-a-Judge, interpretability) $\cdot$ Information Theory.}
\vspace{2pt}
}

% Publications
\compactSection{Publications}{
\begin{itemize}[leftmargin=*,label={},itemsep=4pt]

    \resumeProject{AgriIR: A Scalable Framework for Domain-Specific Knowledge Retrieval}{ECIR 2026}
    {\textbf{Shuvam Banerji Seal}, Aheli Poddar, Alok Mishra, and \href{https://dwaipayanroy.github.io/}{Dr. Dwaipayan Roy}}
    {\item Developed a scalable information retrieval framework specifically designed for agricultural domain knowledge accessibility in developing regions, implementing advanced retrieval algorithms optimized for domain-specific terminology and contextual understanding to advance retrieval systems for social good.
    \item \href{https://doi.org/10.1007/978-3-032-21324-2_37}{\small DOI: 10.1007/978-3-032-21324-2\_37}}

    \resumeProject{Hierarchical Opinion Classification using Large Language Models}{FIRE 2025}
    {\textbf{Shuvam Banerji Seal}, \href{https://alok-mishra6.github.io/}{Alok Mishra}, and \href{https://utkarshaghosh.netlify.app/}{Utkarsha Ghosh}}
    {\item Developed a novel parameter-efficient fine-tuning approach for the Gemma LLM using a custom two-layer classification head and class-weighted cross-entropy loss to address data imbalance. Reformulated three-level hierarchical opinion labels into an 8-class flat scheme, exploring dual methodologies like selective fine-tuning and instruction-tuning to advance LLM adaptation for hierarchical classification under computational constraints.
    \item \href{https://ceur-ws.org/Vol-4173/T10-3.pdf}{\small \url{https://ceur-ws.org/Vol-4173/T10-3.pdf}}}
     % \item \textbf{Technologies:} Python, PyTorch, HuggingFace Transformers, Gemma3-architecture}

    \resumeProject{IISERK@ ToT\_2024: Query Reformulation and Layered Retrieval for Tip-of-Tongue Items}{TREC 2024 Proceedings}
    {\textbf{Shuvam Banerji Seal} \& \textbf{Subinay Adhikary}, Soumyadeep Sar, and \href{https://dwaipayanroy.github.io/}{Dr.\ Dwaipayan Roy}}
    {\item Explored various approaches for known-item retrieval in the TREC 2024 Tip-of-the-Tongue track, focusing on retrieving previously encountered items when searchers struggle to recall exact identifiers
    \item Implemented four-step query reformulation technique combined with two-layer retrieval using BM25 with optimized parameters and Large Language Models for query enhancement
    Achieved best performance with two-layer retrieval approach, obtaining Recall@1000 of 0.8067 and enhanced NDCG metrics through systematic query reformulation
    \item \href{https://trec.nist.gov/pubs/trec33/papers/IISER-K.tot.pdf}{\small \url{https://trec.nist.gov/pubs/trec33/papers/IISER-K.tot.pdf}}}
    % \item \textbf{Technologies:} Python, Lucene (Java), BM25, Large Language Models, Information Retrieval}



    \resumeProject{SycoLex: A Cross-Jurisdictional Benchmark for Detecting Sycophancy in Legal Case Reasoning}{CIKM 2026 Resource Track (under review)}
    {\textbf{Shuvam Banerji Seal} and collaborators, IISER Kolkata}
    {\item Constructed a 1,954-case, three-jurisdiction benchmark (US Supreme Court, Indian Supreme Court, Indian Consumer Courts) with six sycophancy-inducing prompt variants, responses from five LLMs, and paired LLM-as-Judge plus human-expert annotations.
    \item Released openly on Hugging Face under CC-BY-4.0; forms the evaluation resource for the FIRE 2026 SYCO PHANCY shared task.
    \item \href{https://huggingface.co/datasets/ShuvBan/SycoLex}{\small huggingface.co/datasets/ShuvBan/SycoLex}}

    \resumeProject{Computational Modeling of [VO(SALIEP)(DTP)] as Water Reducing Catalyst}{2024-25 (to be submitted)}
    {\textbf{Shuvam Banerji Seal}, Guided by \href{https://soumyajit-roy.com/}{Dr.\ Soumyajit Roy}}
    {\item Implemented DFT methods (B3LYP) for MO energy calculations and electronic structure analysis using Gaussian and conducted mechanistic studies using transition state theory and reaction pathway analysis catalyst performance prediction for water reduction}

    % \resumeProject{On the Nature of Information}{2023-25 (Under Review)}
    % {Guided by \href{https://soumyajit-roy.com/}{Dr.\ Soumyajit Roy}, Dept. of Chemical Sciences, IISER-Kolkata}
    % {\item Developed novel mathematical models for information flow analysis and applied entropy principles to quantify information evolution for a set of constraints to identify and characterize critical inflection points in information dynamics}

\end{itemize}
}

% Master's Thesis
\compactSection{Master's Thesis (Ongoing)}{
\begin{itemize}[leftmargin=*,label={},itemsep=4pt]

    \resumeProject{Project I -- Gray-Code / Karnaugh-Map Structure of Nucleic-Acid and Protein Sequences, Formalised in Lean 4}{2026 -- Present}
    {BS-MS Thesis, IISER Kolkata}
    {\item Encoding DNA, RNA, and protein sequences into binary sequence space and organising the resulting codewords on \textbf{Gray-code} orderings and \textbf{Karnaugh maps}, so that adjacency in the map corresponds to a single-symbol mutation in the biological sequence.
    \item Extracting the binary relations, symmetries, and distance structure induced by this encoding, and testing whether they predict biologically meaningful structure (codon degeneracy, mutation neighbourhoods, amino-acid physicochemical classes).
    \item Formalising and machine-checking every theorem in the \textbf{Lean 4} proof assistant (Mathlib), giving fully rigorous proofs of the structure-based correspondence between the pure combinatorial formulation and real genomic sequences.
    \item \textbf{Keywords:} formal verification, Lean 4, theorem proving, Gray code, Boolean algebra, combinatorics on words, information theory, genomics, bioinformatics.}

    \resumeProject{Project II -- A Transformer Mixture-of-Experts Model for RNA / Nucleic-Acid Structure Prediction}{2026 -- Present}
    {BS-MS Thesis, IISER Kolkata}
    {\item Building a \textbf{Mixture-of-Experts (MoE) transformer} for nucleic-acid (RNA) secondary and tertiary structure prediction, with sparse expert routing so that specialised experts handle distinct structural motifs (helices, loops, pseudoknots) at a fixed inference budget.
    \item Designing the sequence-to-structure pipeline end to end: tokenisation of nucleotide sequences, evolutionary/MSA features, geometry-aware attention, distance-map and torsion-angle prediction heads, and structure-quality losses.
    \item Training and evaluation on public RNA structure benchmarks with multi-GPU distributed training, mixed precision, and systematic ablation of routing strategies and expert counts.
    \item \textbf{Keywords:} deep learning, transformers, mixture-of-experts, RNA structure prediction, structural bioinformatics, PyTorch, distributed training.}

\end{itemize}
}

% Research & Projects
\compactSection{Research Experience}{
\begin{itemize}[leftmargin=*,label={},itemsep=4pt]

    % \resumeProject{Developing an Advanced Retrieval Model for TREC-Tip of Tongue Queries}{2024 (\href{https://trec.nist.gov/pubs/trec33/papers/IISER-K.tot.pdf}{TREC 2024 Proceedings})}
    % {Guided by \href{https://dwaipayanroy.github.io/}{Dr.\ Dwaipayan Roy}, Dept. of Computation and Data-Science, IISER-Kolkata}
    % {Engineered multi-layer BM-25 filtering system (in Lucene, Java) with dynamic search-domain contraction based on extracted keywords and implemented transformer-based query expansion and semantic matching using Local LLMs through multi-shot and chain-of-thought prompting achieving results on par with specifically-modified DPR models}

    % \resumeProject{Hierarchical Opinion Classification with Large Language Models}{2025 (Submitted to FIRE 2025)}
    % {Independent Research with Co-authors: \href{https://alok-mishra6.github.io/}{Alok Mishra} and \href{https://utkarshaghosh.netlify.app/}{Utkarsha Ghosh}}
    % {\item Addressed hierarchical opinion classification using parameter-efficient fine-tuning of Gemma-1B model by attaching a lightweight two-layer classification head and updating only the final transformer block, normalization, and output layers
    % \item Reformulated three-level opinion hierarchy into an 8-class flat scheme and employed class-weighted cross-entropy loss to mitigate data imbalance, ensuring improved treatment of minority categories
    % \item Technology Used: Python, HuggingFace Transformers, PyTorch, Gemma Model}
    \resumeProject{A Hybrid RAG Architecture for Verifiable Answer Synthesis with Automated Source Citation}{July 2025}
    {Under Dr. Dwaipayan Roy, Dept. of Computation and Data-Science, IISER-Kolkata}
    {\item Developed a comprehensive RAG pipeline for IISER-K intranet documents, including automated data acquisition via Selenium, OCR processing with PyMuPDF and EasyOCR, and text extraction
    Implemented context-aware text chunking with metadata enrichment and high-dimensional embeddings using Qwen3-4B model, stored in FAISS vector database
     Designed hybrid search strategy combining dense vector retrieval with BM25 keyword search, fused using Reciprocal Rank Fusion for optimal relevance
    Integrated query refinement with hypothetical answers and verifiable inline citations in a Streamlit application for transparent, grounded responses}
     



     
    % \resumeProject{A Hybrid Retrieval-Augmented Generation Architecture for Verifiable Answer Synthesis with Automated Source Citation}{May 2025 -- July 2025}
    % {Summer Project under Dr. Dwaipayan Roy, Dept. of Computation and Data Science, IISER Kolkata}
    % {\item Engineered a hybrid RAG framework integrating dense (Qwen3-4B embeddings) and sparse (BM25) retrieval fused via Reciprocal Rank Fusion, enabling context-aware document querying, automated source citation, and verifiable answer synthesis through a Streamlit interface for IISER-K intranet resources.
    % \item \textbf{Technologies:} Python, Streamlit, FAISS, LangChain, Qwen3-4B, Gemma3:27B, Selenium, PyMuPDF, EasyOCR, BM25}


    \resumeProject{Developing an Advanced Retrieval Model for TREC-Tip of Tongue Queries}{2024 (\href{https://trec.nist.gov/pubs/trec33/papers/IISER-K.tot.pdf}{TREC 2024 Proceedings})}
    {Guided by \href{https://dwaipayanroy.github.io/}{Dr.\ Dwaipayan Roy}, Dept. of Computation and Data-Science, IISER-Kolkata}
    {\item Engineered multi-layer BM-25 filtering system (in Lucene, Java) with dynamic search-domain contraction based on extracted keywords and implemented transformer-based query expansion and semantic matching using Local LLMs through multi-shot and chain-of-thought prompting achieving results on par with specifically-modified DPR models}
     
\end{itemize}
}




% Industrial Experience
\compactSection{Industrial Experience}{
\begin{itemize}[leftmargin=*,label={},itemsep=4pt]
    \resumeProject{Research Intern, HistoXai (Astroloop Technologies Pvt. Ltd.)}{May 2025 - July 2025}
    {Bangalore}
    {\item Conducted an extensive comparative analysis of digital histopathology slide quality assessment tools, reviewing 30+ open-source frameworks (e.g., HistoQC, PathProfiler, GrandQC, HistoROI, FASTPathology) to evaluate architectural designs, datasets, performance metrics, and limitations for future model integration
    \item \textbf{Technologies:} Python, OpenCV, PyTorch, Scikit-learn, Pandas, Literature Survey Tools}

\end{itemize}
}
% Open Datasets
\compactSection{Open Datasets Released (Hugging Face)}{
\begin{itemize}[leftmargin=*,label={},itemsep=4pt]

    \resumeProject{Hugging Face Profile \href{https://huggingface.co/ShuvBan}{\faDatabase}}{huggingface.co/ShuvBan}
    {Public research datasets (further datasets currently private / pre-release)}
    {\item \textbf{SycoLex} \href{https://huggingface.co/datasets/ShuvBan/SycoLex}{\faDatabase} -- Benchmark for detecting sycophancy in legal case reasoning. 1,954 real legal cases across three jurisdictions (US Supreme Court 300, Indian Supreme Court 1,500, Indian Consumer Courts 154), six sycophancy-inducing prompt variants, responses from five LLMs, LLM-as-Judge and human expert annotations. Tasks: text classification, text generation. License CC-BY-4.0.
    \item \textbf{AgriIR\_dataset} \href{https://huggingface.co/datasets/ShuvBan/AgriIR_dataset}{\faDatabase} -- 15,247 curated Indian agricultural knowledge entries with quality scores and domain metadata (crops, soils, regions, farming practices), released alongside the AgriIR paper (ECIR 2026; \href{https://arxiv.org/abs/2604.16353}{arXiv:2604.16353}). License CC-BY-4.0.
    \item Datasets are versioned, documented with dataset cards, and used as the public evaluation resources for the SYCO PHANCY (FIRE 2026) shared task and the AgriIR retrieval framework.}

\end{itemize}
}

% Entrepreneurial Ventures
\compactSection{Entrepreneurial Ventures}{
\begin{itemize}[leftmargin=*,label={},itemsep=4pt]

    \resumeProject{Co-Founder \& Lead Developer, iFiNN \href{https://synapse-iiserk.github.io/}{\faGlobe}}{2025 -- Present}
    {Incubated at RISE Foundation, IISER Kolkata}
    {\item Co-founded an AI-driven fintech startup funded by \textbf{MeitY Startup Hub (MSH)} through the \textbf{GENESIS} program.
    \item Developing a user-friendly platform with simple UI/UX to democratize access to essential ML and AI tools for automated smart trading alerts across stocks, crypto, and mutual funds.
    \item \textbf{Status:} Funded; DPIIT registration in progress (Q1 2026); Product shipping by 2027.}

    \resumeProject{Co-Founder \& CTO, UnderWater AI \href{https://underwater-ai.github.io/}{\faGlobe}}{2025 -- Present}
    {Incubated at RISE Foundation, IISER Kolkata}
    {\item Co-founded a deeptech startup funded by \textbf{MeitY Startup Hub (MSH)} through the \textbf{GENESIS} program.
    \item Developing software solutions to enhance underwater image quality using deep neural networks and early fusion learning for marine species identification, serving researchers and industrial marine applications.
    \item \textbf{Status:} Funded; DPIIT registration in progress (Q1 2026); Product shipping by 2027.}

\end{itemize}
}

\compactSection{Flagship Open-Source Software \& Research Frameworks}{
\begin{itemize}[leftmargin=*,label={},itemsep=4pt]

    \resumeProject{Molecule3D / LAMMPS Web GUI -- Browser-Native Molecular Simulation Workbench \href{https://github.com/Shuvam-Banerji-Seal/lammps-web-gui}{\faGithub}\ \href{https://shuvam-banerji-seal.github.io/lammps-web-gui/}{\faGlobe}}{2026 -- Present}
    {Minimum viable product; free, open-source (MIT), GPU-accelerated, runs entirely client-side}
    {\item Built a complete \textbf{LAMMPS} workbench in the browser that lowers the entry barrier to molecular dynamics: a visual \textbf{Script Builder} (drag-and-drop flowchart over a 186-command library covering the full general-command surface, locked by an automated coverage test against \texttt{docs.lammps.org}), a \textbf{Compiler Helper} that turns 27 CMake build options (Kokkos, GPU back end/precision/arch, FFT/heFFTe, OpenMP, FFMPEG) into ready-to-run LAMMPS build scripts, and a 3D \textbf{Structure Viewer}.
    \item Reads LAMMPS \texttt{data}, XYZ, PDB and CIF structures plus LAMMPS \texttt{dump} trajectories; supports trajectory playback, in-browser MP4/WebM recording, distance/angle/dihedral measurement, and an \textbf{Analysis} module computing PBC-aware RDF $g(r)$, MSD, density profiles and speed histograms with CSV export.
    \item Round-trips existing \texttt{in.*} input scripts back into an editable flowchart, exports presentation-quality SVG/PNG pipeline diagrams, and persists all state locally -- nothing is uploaded.
    \item \textbf{Stack:} TypeScript, React, Three.js / React Three Fiber, WebGL, Vite, Vitest (143 passing tests), GitHub Actions CI/CD.}

    \resumeProject{AgriIR -- Configurable RAG Framework for Domain-Specific Retrieval \href{https://github.com/Shuvam-Banerji-Seal/AgriIR}{\faGithub}}{2025 -- 2026}
    {Reference implementation of the ECIR 2026 paper}
    {\item Six-stage configurable retrieval-augmented generation pipeline with deterministic citations and domain-specific agents, built on FAISS, Ollama-served open-weight LLMs, and hybrid BM25 + dense retrieval.
    \item Ships the accompanying 15k-entry open dataset and a reproducible evaluation harness. \href{https://doi.org/10.1007/978-3-032-21324-2_37}{\small DOI: 10.1007/978-3-032-21324-2\_37}}

    \resumeProject{IISERKonnect -- Campus Super-App for IISER Kolkata \href{https://github.com/Shuvam-Banerji-Seal/iiserkonnect-website}{\faGithub}\ \href{https://shuvam-banerji-seal.github.io/iiserkonnect-website/}{\faGlobe}}{2026 -- Present}
    {Sole developer; native Android (Kotlin) + web client. \textbf{In campus-wide testing with 2,000+ users at IISER Kolkata}}
    {\item Designed and built a single application that unifies the entire IISER Kolkata intranet surface -- WeLearn courses, files, assignment deadlines and attendance; mess transactions, budget analytics and daily menu; academic calendar and campus events; grade cards; the previous-year-question archive; student and administrative notice boards; the library catalogue with live availability; the ePrints research archive; the VoIP directory; the TCP counter and MAC registration; and a network health dashboard.
    \item Implemented every campus service as a hand-written HTML parser (Kotlin/Jsoup on Android, ported 1:1 to \texttt{DOMParser} on web), with captcha-backed login flows, per-session cookie jars, and a strict campus-only host allow-list; credentials never leave the device.
    \item Added a serverless \textbf{Campus Chat} over WebRTC data channels (DTLS end-to-end encrypted, invite-code handshake, optional STUN for strict NATs), replacing the earlier raw-TCP discovery layer.
    \item \textbf{Stack:} Kotlin, Jetpack Compose, Jsoup, OkHttp, JavaScript (ES modules), Node.js companion proxy, WebRTC, Playwright end-to-end tests.}

    \resumeProject{Introduction to Molecular Simulation -- Open Talk \& Hands-On Companion \href{https://github.com/Shuvam-Banerji-Seal/introduction-to-molecular-simulation}{\faGithub}}{2026}
    {Most-starred public repository; Marp slide deck, concept notes and runnable examples}
    {\item Authored a complete beginner-to-running-simulation course: statistical-mechanics motivation, force fields, MD integrators/thermostats/barostats vs.\ Monte Carlo, NVE/NVT/NPT ensembles, and observable extraction (RDF, MSD, VACF, free energy), with reproducible LAMMPS simulations used as the demo trajectories for Molecule3D.}

    \resumeProject{DFT Notes -- Open Knowledge Base on Density Functional Theory \href{https://github.com/Shuvam-Banerji-Seal/DFT-notes}{\faGithub}\ \href{https://shuvam-banerji-seal.github.io/DFT-notes/}{\faGlobe}}{2026}
    {Open educational resource}
    {\item Progressive chapters from the Schr\"odinger and many-body problem through Hartree-Fock, the Hohenberg-Kohn theorems, the Kohn-Sham equations and exchange-correlation functionals, each with runnable Python and generated plots, plus a Mermaid-driven knowledge graph of the chapter dependencies.}

    \resumeProject{Fernholz Stochastic Portfolio Theory Python Library \href{https://github.com/XAheli/Fernholz-SPT}{\faGithub}}{Ongoing}
    {Quantitative Finance Library}
    {\item Developing a Python library for Stochastic Portfolio Theory, implementing diversity-weighted portfolios and relative arbitrage strategies using stochastic calculus and continuous semimartingales.}

\end{itemize}
}

\compactSection{Tutorials Developed}{
\begin{itemize}[leftmargin=*,label={},itemsep=4pt]
    \resumeProject{Complete C Programming Course \href{https://github.com/Shuvam-Banerji-Seal/C-Programming-for-Beginners}{\faGithub}}{2025}
    {Open Source Curriculum}
    {\item Created a comprehensive 20-module C programming curriculum covering fundamentals to advanced topics including Network Programming, Machine Learning, and GUI development with GTK4.}

    \resumeProject{Python Course for Beginners \href{https://github.com/Shuvam-Banerji-Seal/Python-Course-for-Beginners}{\faGithub}}{2025}
    {Open Source Curriculum}
    {\item Designed a progressive Python learning resource spanning core concepts to local LLM deployment, featuring interactive notebooks, real-world project implementations, and database integration.}

    \resumeProject{Shaders for Beginners \href{https://github.com/Shuvam-Banerji-Seal/shaders-for-beginners}{\faGithub}\ \href{https://shuvam-banerji-seal.github.io/shaders-for-beginners/}{\faGlobe}}{2026}
    {Open Source Curriculum with interactive WebGL editor}
    {\item Wrote 11 progressive lessons teaching GPU shader programming from first principles -- OpenGL and GLSL, the graphics pipeline, an introduction to Vulkan, and audio-reactive shaders -- with 15+ worked shaders and C/C++ host implementations.}

    \resumeProject{CH4114 Molecular Simulation -- Hands-On Tutorials \href{https://github.com/Shuvam-Banerji-Seal/CH4114_Molecular_Simulation_Tutorials}{\faGithub}}{2026}
    {Course tutorial material, IISER Kolkata}
    {\item Prepared and delivered tutorial sessions on reproducible computational practice for molecular simulation: Git and GitHub, project organisation, and \texttt{uv}-based Python environment management.}
\end{itemize}
}

\compactSection{Projects}{
\setlength{\columnsep}{15pt}
\begin{multicols}{2}
\begin{itemize}[leftmargin=*,label={},itemsep=4pt]

    \resumeProject{Solvation of C$_{60}$ Nanoparticles in Water \href{https://github.com/Shuvam-Banerji-Seal/Solvent-Structure-around-NanoParticles}{\faGithub}}{2025}
    {Computational Science Project}
    {\item Ran LAMMPS molecular dynamics of C$_{60}$ solvation with AIREBO potentials, automated structural feature extraction (RDF, coordination numbers), and trained a surrogate ML pipeline (Gaussian process regression, neural networks, XGBoost) to predict system properties at unsampled parameter values.}

    \resumeProject{opngx -- High-Speed Camera Footage Extractor \href{https://github.com/Shuvam-Banerji-Seal/opngx}{\faGithub}}{2026}
    {Scientific Instrumentation Tooling}
    {\item Wrote a pixel-exact extractor for Optronis \texttt{.bin} high-speed-camera recordings that saturates all CPU cores with runtime SIMD dispatch, shipped as a CLI, a Tk GUI and a Python library for Linux and Windows.}

    \resumeProject{TCP Activity Monitor for Arch Linux \href{https://github.com/Shuvam-Banerji-Seal/TCP-Count-Monitor-for-Arch-Linux}{\faGithub}}{2025}
    {System Utility}
    {\item Developed a system-wide TCP monitoring utility for Arch Linux integrating systemd services to log kernel statistics, socket states, and per-process connections, enhancing network diagnostic capabilities.}

    \resumeProject{WeLearn Bot in C \href{https://github.com/Shuvam-Banerji-Seal/welearnbot_in_C}{\faGithub}}{2025}
    {Utility Application}
    {\item Engineered a multi-threaded educational resource automation tool in C featuring both CLI and GTK4 GUI interfaces with secure session management and encrypted credential storage for efficient bulk downloads.}

    \resumeProject{IndicAgri: AI Platform for Indian Agriculture \href{https://github.com/Shuvam-Banerji-Seal/Answering_Agriculture}{\faGithub}}{Aug 2025}
    {Capital One Launchpad Hackathon 2025}
    {\item Architected a multi-modal RAG chatbot for Indian agriculture supporting 20+ languages using Gemma/DeepSeek LLMs and Indic-Conformer, backed by an autonomous 15k+ document curation pipeline.}
    
    % \resumeProject{Automated WhatsApp Bulk Message Sender \href{https://github.com/Shuvam-Banerji-Seal/Automated-Whatsapp-Bulk-Message-Sender-using-Google-Forms-Sheets-and-Python}{\faGithub}}{2025}
    % {Independent Project}
    % {\item Created an automated bulk messaging bot using Python and Selenium with Google Sheets API integration, ensuring robust error handling, Unicode support, and seamless session persistence.}
     
    \resumeProject{Automated Event Coupon Management System \href{https://github.com/Shuvam-Banerji-Seal/Automated-Event-Coupon-Sender-Email-and-Verification-Application}{\faGithub}}{2025}
    {Event Automation System}
    {\item Built a secure Flask-based event management system featuring Google OAuth 2.0 authentication, automated bulk emailing, and real-time QR code verification with mobile scanner integration.}
     
    \resumeProject{Legal Document Retrieval RAG Application \href{https://github.com/Shuvam-Banerji-Seal/CCA-2015-LLM}{\faGithub}}{Feb 2025}
    {Legal Tech RAG System}
    {\item Developed a legal document retrieval RAG stack leveraging LLaMA-3.2, Nomic embeddings, and ChromaDB, with a Streamlit interface and Mistral-OCR for processing complex legal texts.}
     
    \resumeProject{Full-Stack E-Commerce GUI in C \href{https://github.com/Shuvam-Banerji-Seal/CS3101-E-Commerce-App-in-C.git}{\faGithub}}{2024}
    {Under \href{https://www.iiserkol.ac.in/web/en/people/faculty/cds/kripaghosh/}{Dr.\ Kripabandhu Ghosh}, IISER-Kolkata}
    {\item Implemented a full-stack E-Commerce GUI in C using GTK4 and SQLite3, incorporating a static context chatbot with optimized BM25 retrieval for intelligent customer interaction.}
     
    \resumeProject{Wi-Fi Channel Optimizer \href{https://github.com/Shuvam-Banerji-Seal/Wi-Fi-channel-optimizer}{\faGithub}}{2024}
    {Network Optimization Tool}
    {\item Engineered a Bash-based Wi-Fi optimization script to automate network scanning, speed benchmarking, and channel selection, significantly enhancing wireless network performance on Linux systems.}
     
    \resumeProject{SMC Canteen System Modernization}{2024}
    {IISER-Kolkata}
    {\item Modernized the IISER-Kolkata Canteen System by porting legacy code to a modular Django-based architecture, improving maintainability and scalability for 2000+ users.}

    \resumeProject{Coordination Chemistry Twists Simulator \href{https://github.com/Shuvam-Banerji-Seal/Ray-Dutt-and-Bailar-Twists-Simulator}{\faGithub}}{Nov 2024}
    {Computational Chemistry Simulation}
    {\item Developed a Python molecular dynamics simulator for visualizing Ray-Dutt and Bailar twist mechanisms in coordination complexes, providing interactive 3D stereochemical analysis.}
     
    \resumeProject{Agentic Database Builder \href{https://github.com/Shuvam-Banerji-Seal/Agentic_Database_Builder}{\faGithub}}{Aug 2024}
    {Autonomous Data Engineering}
    {\item Designed an AI-driven agentic system for autonomous database construction, utilizing LLM decision-making for automated schema generation, data validation, and scalable curation.}
     
    \resumeProject{Interactive Spherical Harmonics Visualizer \href{https://github.com/Shuvam-Banerji-Seal/simple-spherical-harmonics-visualizer}{\faGithub}}{2023-24}
    {Scientific Visualization Engine}
    {\item Built an interactive 3D quantum orbital visualizer using Flask and Plotly, enabling real-time computation and rendering of spherical harmonics for educational visualization.}
     
\end{itemize}
\end{multicols}
\vspace{-5pt}
\begin{center}
    \small{\textit{More projects can be found at my \href{https://shuvam-banerji-seal.github.io}{Portfolio Website} or \href{https://github.com/Shuvam-Banerji-Seal}{Github Profile}.}}
\end{center}
}






% Technical Skills
\compactSection{Technical Skills}{
\begin{itemize}[leftmargin=0.15in,label={},itemsep=2pt]
    \skillBlock{Core Competencies}{Technical Leadership | Research \& Development | Algorithm Development | Scientific Software Engineering | Strategic Planning | Event Management | Mentoring | Public Speaking | Cross-functional Collaboration}
    \skillBlock{Research Computing}{Information Retrieval (Apache Lucene, BM25, Reciprocal Rank Fusion) | High-Performance Computing (HPC) | Parallel Computing (OpenMP, MPI, CUDA) | Numerical Methods | Statistical Mechanics | Information Theory | Reproducible Research}
    \skillBlock{Computational Chemistry \& Materials}{Density Functional Theory (DFT) | \textbf{Quantum ESPRESSO} (plane-wave/pseudopotential SCF, geometry relaxation, band structure, DOS, phonons) | Gaussian (B3LYP, basis sets, transition-state search) | VASP workflows | ORCA | Molecular Dynamics (LAMMPS, GROMACS, VMD, OVITO) | Force Fields (AIREBO, ReaxFF, EAM, Lennard-Jones) | Monte Carlo | Ab-Initio Electronic Structure | Exchange-Correlation Functionals | Radial Distribution Functions, MSD, Free-Energy Analysis}
    \skillBlock{Structural Biology \& Bio-Informatics}{Protein and Nucleic-Acid Structure Prediction | RNA Secondary/Tertiary Structure | Sequence Analysis (DNA, RNA, Protein) | PyMol | ChimeraX | AutoDock Vina | PyDock | Molecular Docking | Genomic Data Analysis}
    \skillBlock{Machine Learning \& Deep Learning}{PyTorch | TensorFlow | Transformers | Mixture-of-Experts (MoE) | Attention Mechanisms | Parameter-Efficient Fine-Tuning (LoRA/QLoRA) | Distributed \& Mixed-Precision Training | Gaussian Process Regression | XGBoost | Scikit-learn | Model Evaluation \& Ablation Studies}
    \skillBlock{LLMs, RAG \& Agentic AI}{Retrieval-Augmented Generation (RAG) | Dense \& Sparse Hybrid Retrieval | Vector Databases (FAISS, ChromaDB) | LangChain | Ollama | vLLM | HuggingFace Hub \& Datasets | Prompt Engineering (multi-shot, chain-of-thought) | Autonomous Agents | LLM-as-a-Judge | Sycophancy \& Bias Evaluation | AI Safety | Qwen2/3 | Gemma (1B, 27B) | DeepSeek | LLaMA-3.2 | Nomic Embeddings | Mistral-OCR | Indic-Conformer}
    \skillBlock{Formal Methods}{Lean 4 \& Mathlib | Machine-Checked Theorem Proving | Formal Verification | Boolean Algebra \& Gray Codes | Combinatorics on Words}
    \skillBlock{Programming Languages}{Python | C/C++ | Kotlin | TypeScript/JavaScript | Java | Rust | Fortran | Lean 4 | Bash | SQL | QBASIC | GWBASIC}
    \skillBlock{Python Libraries - Core \& Scientific}{NumPy | Pandas | SciPy | Matplotlib | Plotly | Scikit-learn | OpenCV | PyTorch | TensorFlow | ASE | MDAnalysis}
    \skillBlock{Python Libraries - NLP, Data \& Documents}{HuggingFace Transformers \& Datasets | LangChain | NLTK | SpaCy | FAISS | ChromaDB | Selenium | BeautifulSoup | PyMuPDF | EasyOCR | Streamlit | Manim}
    \skillBlock{Digital Pathology Tools}{HistoQC | PathProfiler | GrandQC | HistoROI | FASTPathology}
    \skillBlock{Development Stack}{Full Stack (HTML/CSS/JavaScript) | React | Three.js / React Three Fiber | WebGL \& GLSL | Vite | Django | Flask | Node.js | REST APIs | WebRTC | Databases (PostgreSQL, MongoDB, MySQL, SQLite3) | Git \& GitHub Actions (CI/CD) | Docker | CMake | Make | Vitest | Playwright}
    \skillBlock{Mobile Development}{Android (Kotlin) | Jetpack Compose | Coroutines \& Flow | Room | OkHttp | Jsoup}
    \skillBlock{GUI Development}{GTK4 in C | Qt in C++ | Glade | Tkinter}
    \skillBlock{Operating Systems \& Infrastructure}{Linux (Arch, Ubuntu) | Shell \& systemd | SSH, OpenSSL | Server \& Cluster Handling | SLURM-style Job Submission | GPU Computing (CUDA)}
    \skillBlock{Laboratory Instrumentation}{UV-Vis | ATR-FTIR | TGA | DSC | Optical Bench | XRD | GC-MS | Column Chromatography | Fluorimeter | SEM | TEM | AFM}
    \skillBlock{Additional Tools}{\LaTeX\ \& Beamer | Marp | UI/UX Design (Figma) | Technical Writing | Scientific Visualisation | HTML E-mailing}
\end{itemize}
}




% Achievements
\compactSection{Achievements \& Honors}{
\begin{itemize}[leftmargin=*,label={},itemsep=4pt]
    \large \item \textbf{Hackathons:}  
    \resumeSubheading{\href{https://capitalone.hackerearth.com/}{Capital One Launchpad} (Top 14 Team among 5000+ Teams) \href{https://github.com/Shuvam-Banerji-Seal/Answering_Agriculture}{\faGithub}}{2025}
    {Developed IndicAgri: Multi-modal RAG Platform for Indian Agriculture (Team Fibonacci)}{Capital One India}
    \denseList{\small{
        \item Selected as one of the Top 14 teams from 5,073 registered teams across India in Capital One's flagship AI/ML innovation challenge \vspace{-2pt}
        \item Built an agentic AI-powered agricultural advisor supporting 20+ Indian languages with scientifically cited, hyper-localized outputs using RAG, LangChain, and open-source LLMs (Gemma, DeepSeek) \vspace{-2pt}
        \item Engineered autonomous data pipeline creating a novel 15k+ document dataset for Indian agriculture, released on Hugging Face
    }} \vspace{2pt}
    \resumeSubheading{\href{https://statuscode-1.devfolio.co/}{StatusCode1} (Awarded 1st Rank in GIAN Track) \href{https://github.com/Shuvam-Banerji-Seal/LLM-based-Searcher-for-GIAN-s-Abandoned-US-Patents}{\faGithub}}{2024}
    {Developed an AI-based Search Engine for \href{https://gian.org/abandoned-us-patents/}{GIAN's Database For Abandoned US Patents}}{IIIT-Kalyani}
    \denseList{\small{
        \item A searching algorithm based on Nomic Embeddings of the patent abstracts and similarity computations thus enabling the user to search for patents in natural language without the use of specific terminologies. \vspace{-2pt}
        \item Also created a web-scraping algorithm to get the patent data using Selenium, BeautifulSoup and fake\_useragent for anti-scraping measures by the website
    }} \vspace{2pt}
    % \resumeSubheading{\href{https://www.sih.gov.in/}{Smart India Hackathon} (Participated)}{2024}
    % {Problem ID: SIH1701 and SIH1734}{SIH-2024}
    % \denseList{\small{
    %     \item Proposed a "Context-Aware AI Framework for Commerical Courts (2015) Cases' Retrieval and Predictive Analytics" utilizing RAG, based on embeddings produced from the Legal documents (1701) \vspace{-2pt}
    %     \item Proposed a "Next-Gen Air Quality Mapping: A Deep Learning Approach with SRDIFF and Kolmogorov Arnold Networks" (1734)
    % }} \vspace{2pt}
    
        \resumeSubheading{\href{https://statuscode0.devfolio.co/}{StatusCode0} (Awarded 1st Rank in MATLAB Track)}{2023}
    {Developed a Domestic Waste Type Data analysis tool for a proposed Start-Up Solution}{IIIT-Kalyani}
 
 \vspace{4pt}
    \large \item \textbf{National Level Basic Sciences Competitions:}
        \resumeSubheading{ChemEnigma (1st Rank)}{2025}
    {Emerged champions at 72 hours Chemistry contest of theoretical, experimental and concept-presentation}{IISc-Bangalore}
    \resumeSubheading{All Bengal Chemistry Quiz (2nd Runners Up)}{2025}
    {Solved chemistry based questionnaires under time constraints}{Presidency University}
    \resumeSubheading{Mimansa (Zonal Topper)}{2024}
    {Contributed to the team in Mathematical problem solving}{IISER-Pune}
   \resumeSubheading{NAEST-National Anveshika Experimental Skill Test (Zonal Runners Up)}{2023}
    {Create extensive experimental setup using homely items- see \href{https://github.com/Shuvam-Banerji-Seal/NAEST_Sample_Experiments}{papers}}{NANI, IIT Kanpur and Shiksha Sopan}  


    \item \textbf{Competitive Examinations:}
    \resumeSubheading{Qualified JEE Mains and Advanced}{2022}
    {Ranked in Top 0.1 fraction of Candidates}{}\vspace{-4pt}
    \resumeSubheading{Qualified IAT(IISER-Aptitude Test)}{2022}
    {Ranked in top 0.06 fraction of candidates}{}\vspace{-4pt}
    \resumeSubheading{Qualified WBJEE}{2022}
    {Ranked in Top 0.05 fraction of Candidates}

    

    \vspace{4pt}\item \textbf{Scholarships and Honors:}
    \resumeSubheading{Reliance Foundation Undergraduate Scholar}{2023}
    {Qualified the RF-UG Aptitude test to be in the top 5000 students to be awarded this} {}\vspace{-2pt}
    \resumeSubheading{Best Young Scientist Speaker on Nanotechnology}{2019}
    {Successfully presented at the prestigious conference }{World Science Conference, Jadavpur University}
    
\end{itemize}
}

%     \item \textbf{Hackathons:}
%     \denseList{
%         \item 1st Position in \href{https://statuscode-1.devfolio.co/}{StatusCode1} (GIAN track) organized by IIIT-Kalyani - AI Search Engine \& Patent Analysis Tool
%         \item Reliance Foundation Undergraduate Scholar (2023)
%         \item Most Creative Use of MatLab - StatusCode0 Hackathon, IIIT Kalyani
%         \item NAEST Zonal Runners Up (2023)
%     }
%     \item \textbf{Academic Excellence:}
%     \denseList{
%         \item JEE Advanced (Top 0.1\%), WBJEE (Top 0.05\%), IAT (Top 0.06\%) - 2022
%         \item Zonal Topper - Mimansa Mathematics Competition, IISER Pune (2024)
%         \item SOF IEO: International Rank 331 (2020-21), 536 (2019-20)
%         \item SOF NSQ: International Rank 346, Zonal Rank 22 (2020-21)
%     }
% \end{itemize}

% Experience
\compactSection{Professional Experience}{
\begin{itemize}[leftmargin=*,label={},itemsep=4pt]

    \resumeProject{Freelance Web Developer -- ChemActiva Innovations Pvt.\ Ltd.\ \href{https://chemactiva.com/}{\faGlobe}}{2024 -- 2025}
    {Nanocellulose deeptech startup; DST-NIDHI PRAYAS and HDFC Parivartan grantee}
    {\item Designed and built the company's public website end to end as a freelance engagement, covering information architecture, responsive layout, product and technology pages, and deployment.
    \item \textbf{Live at} \href{https://chemactiva.com/}{chemactiva.com}.}

    \resumeProject{Web Development for Anicon 3.0 \href{https://anicon3.github.io/}{\faGlobe}}{2024-2025}
    {Inquivesta XI, IISER Kolkata}
    {\item Developed and led the web development of Anicon 3.0 Event. }

    \resumeProject{Web Development for Material Science Laboratory \href{https://shuvam-banerji-seal.github.io/EFAML_WEB/index.html}{\faGlobe}}{2025}
    {EFAML, IISER Kolkata}
    {\item Developed and designed the researchers lab info page available under Dr. Soumyajit Roy's homepage.}
     
    \resumeProject{Private Educator \& Technical Trainer}{2018-Present}
    {Advanced Computing \& Basic Sciences Instruction \textit{(Self-Employed)}, Kolkata}
    {\item Developed and conducted courses in Computer Science, Physics, Chemistry and English for High School students (ICSE, CBSE, WB Board)
    \item Mentored 50+ students for Board and competitive examinations}
     
    \resumeProject{Technical Consultant}{2021-Present}
    {Self-Employed, Kolkata}
    {\item Designed and implemented high-performance computing solutions
    \item Expertise in system optimization, BIOS/UEFI configuration, and OS installation
    \item Successfully completed 50+ custom build projects with 100\% client satisfaction}
     
    \resumeProject{Published Author}{2020}
    {MindScapes (ISBN: 978-9389923209), Kolkata}
    {\item Published creative anthology work focusing on metaphorical and philosophical themes
    \item Conducted workshops on technical and creative writing}
     
\end{itemize}
}



% Leadership
\compactSection{Leadership \& Community Impact}{
\begin{itemize}[leftmargin=*,label={},itemsep=2pt]
    \resumeProject{President, Slashdot - The Programming and Design Club \href{https://slashdot-iiserk.github.io/}{\faGlobe}}{Aug 2025 -- Present}
    {IISER Kolkata -- \textbf{elected with a 99\% approval vote of the club membership}}
    {\item Elected to lead the official programming and design club of IISER Kolkata with a \textbf{99\% mandate}, one of the strongest endorsements recorded for the post, reflecting sustained contribution to the club before taking office.
    \item Directing the club's technical programme: workshops, hackathons, machine-learning and local-LLM sessions, and the club's open-source and web presence, building a durable coding culture on campus.}

    \resumeProject{Office Bearer, Valence - The Chemistry Club}{Aug 2025 -- Present}
    {IISER Kolkata}
    {\item Leading the official chemistry club of IISER Kolkata, organizing seminars, student talks, and departmental events to foster a vibrant chemical science community.}

    \resumeProject{Organizer, Qiskit Fallfest 2025}{Oct 2025}
    {Sponsored by IBM Quantum}
    {\item Organized the local chapter of Qiskit Fallfest (one of 150 selected institutes globally sponsored by IBM); conducted hands-on workshops and tutorials on Qiskit SDK and quantum computing fundamentals.}

    \resumeProject{Event Management \& Logistics}{2023 -- 2025}
    {IISER Kolkata}
    {\item \textbf{Anicon 3.0 \& 2.0:} Led organization of two editions of the sci-fi convention with 500+ participants.
     \item \textbf{Academic Events:} Organized Supra-Molecular Discussions 2024 and GIAN Courses on Soft-Oxometalates and X-Ray Crystallography.}

    \resumeProject{Social Impact \& Community Service}{2020 -- 2021}
    {Volunteering}
    {\item Served as a COVID-19 relief coordinator and currently mentoring students under the Ek-Pehal education initiative.}

    \item \textbf{Additional Activities:} \small {District Level Debate and Quiz finalist | Shotokan Karate practitioner | 4th Year Art and Painting Student (Stroke Art \& Portraits)}
\end{itemize}
}


\compactSection{Honors \& Awards}{
\begin{itemize}[leftmargin=*,label={},itemsep=4pt]

    \resumeProject{\textbf{1st Prize (Rs.~2,00,000) -- UIDAI Data Hackathon 2026} \href{https://event.data.gov.in/challenge/uidai-data-hackathon-2026/}{\faGlobe}}{Jan--May 2026}
    {Unique Identification Authority of India (UIDAI), NIC and MeitY, Government of India}
    {\item Won the \textbf{national Aadhaar data hackathon} outright: nearly \textbf{15,000 teams registered}, over \textbf{5,000 submitted solutions}, 30 projects were shortlisted and 15 finalist teams were assessed by the jury before the winning entry was announced at the valedictory session on \textbf{9 May 2026}.
    \item Winning team drawn from \textbf{IISER Kolkata} and the Institute of Engineering and Management, Kolkata. Delivered an in-depth analysis of UIDAI's aggregated Aadhaar enrolment and update datasets, characterising \textbf{biometric update patterns and enrolment trends across regions, states and demographic groups} and deriving predictive indicators.
    \item Produced concrete operational recommendations for improving Aadhaar service delivery; the CEO of UIDAI publicly noted that such work has ``the potential to directly support policy and operational improvements''.
    \item \textbf{Coverage:} \href{https://www.pib.gov.in/PressReleasePage.aspx?PRID=2259168}{\small Press Information Bureau} $\mid$ \href{https://uidai.gov.in/en/media-resources/media/aadhaar-in-prints/19821-uidai-data-hackathon-2026-showcases-data-driven-innovations-for-inclusive-governance.html}{\small UIDAI} $\mid$ \href{https://www.tribuneindia.com/news/business/uidai-data-hackathon-2026-highlights-student-led-digital-identity-innovations/}{\small The Tribune}
    \item \textbf{Keywords:} large-scale data analysis, population-scale identity data, exploratory data analysis, geospatial and demographic segmentation, predictive indicators, public-policy analytics, data visualisation, Python, Pandas.}

\end{itemize}
}

\compactSection{Invited Talks and Positions of Jury}{
\begin{itemize}[leftmargin=*,label={},itemsep=4pt]



    \resumeProject{Technical Talk – Running and Optimizing Local LLMs}{Aug 2025}
    {Slashdot Student Chapter, IISER Kolkata}
    {
        \item Delivered a hands-on session covering hardware setups, quantization methods, embedding pipelines, vector databases, and practical optimization techniques for local LLM deployment.
    }

    \resumeProject{Invited Talk – Sustainability, AI, and Emerging Technologies}{Aug 2025}
    {Valence – Chemistry Society, IISER Kolkata}
    {
        \item Presented an interdisciplinary talk bridging sustainability science with AI-driven modelling, retrieval systems, and computational chemistry tools.

    }
    \resumeProject{Track Co-organizer, FIRE 2026 -- SYCO PHANCY Shared Task \href{https://sycolex.com/}{\faGlobe}}{2026}
    {Explainable AI in Legal Reasoning: From Statute Prediction to Sycophancy Detection}
    {\item Co-organizing a shared task track at FIRE 2026 with international collaborators from IISER Kolkata, Université de Bretagne Occidentale (France), and University of Amsterdam (Netherlands), evaluating LLM interpretability and sycophancy detection in legal AI.}

    \resumeProject{Guest Lecturer \& Guest of Honour -- CBSE District-Level Deliberation (DLD), Kolkata District \href{https://shuvam-banerji-seal.github.io/ai-in-real-world-context/}{\faGlobe}\ \href{https://github.com/Shuvam-Banerji-Seal/ai-in-real-world-context}{\faGithub}}{Aug 2026}
    {CBSE STEM DLD workshop on Computational Thinking \& AI, Abhinav Bharati High School (ABHS), Kolkata}
    {\item Invited by the Central Board of Secondary Education as an \textbf{AI subject expert} to address \textbf{64 teachers from 38 schools} across the Kolkata district on ``AI in the Real-World Context'' -- what today's AI tools and services actually are, the next-token-prediction machinery underneath them, and how to teach with and about them responsibly.
    \item Authored and released the full 60-minute talk openly in two formats: an interactive, pointer-driven \href{https://shuvam-banerji-seal.github.io/ai-in-real-world-context/}{HTML slide deck} (26 slides, keyboard navigation, in-code SVG diagrams) and a \href{https://shuvam-banerji-seal.github.io/ai-in-real-world-context/main.pdf}{LaTeX Beamer PDF} with matching TikZ figures, under a public GitHub repository so attending schools can reuse the material.}

    \resumeProject{Invited Judge -- Voice of Youth 2026 National Summit \& Hackathon}{1--2 Aug 2026}
    {Raajkutir, Kolkata; national teen-led innovation and leadership summit supported by Mr.\ Harshavardhan Neotia}
    {\item Served on the judging panel evaluating shortlisted student teams presenting at the National Offline Summit across the technology and innovation domains, assessing technical depth, feasibility, and demonstrated impact.}

    \resumeProject{Domain Judge – Computer Science Events}{July 2025}
    {La Martiniere for Boys Annual Science Fest, Kolkata}
    {
        \item Invited to serve as an external judge for multiple Computer Science events involving project demonstrations, live coding, and problem-solving rounds.
    }
\end{itemize}
}

% Education (Moved to End)
\compactSection{Education}{
\begin{itemize}[leftmargin=*,label={},itemsep=3pt]
    \resumeSubheading{Indian Institute of Science Education and Research - Kolkata}{CGPA: 8.2}
    {BS-MS (Chemistry Major, Computer Science Minor)}{2022-2027 \textit{(expected)}}
    \vspace{-5pt}
    \resumeSubheading{Calcutta University}{CGPA: 8.308}
    {B.Sc Honours in Physics (1st Year Only)}{2021-2022}
        \vspace{-5pt}
    \resumeSubheading{Jodhpur Park Boys' High School}{83\%}
    {Higher Secondary Education in Physics, Mathematics, Chemistry, Computer Science (WBCHSE)}{2019-2021}
        \vspace{-5pt}
    \resumeSubheading{The New Horizon High School}{83.75\%}
      {Secondary Level Schooling(English Medium) under WBBSE}{2009-2019}
\end{itemize}
}


\vspace{-0.4cm}
% Add referrals line at the bottom right
\vspace*{\fill}
\begin{flushright}
\textit{References available upon request}
\end{flushright}

\end{document}